%% Stone representation, reconstruction, and functorial duality.

\subsection{Ultrafilters and Stone representation}
\label{sec:p10-stone-representation}

The proofs in this subsection work in \(\mathsf{ZF}\) together with
the Boolean Prime Ideal Theorem (BPIT), in its proper-filter-extension
form stated above. This assumption suffices for the arbitrary Boolean
algebras considered here. The complete source presentation is
\cite[IV\S\S3--4]{BurrisSankappanavar2012}.

\begin{hblemma}{thm:p10-ultrafilter-separation}{proved=annex:p10.stone.1}{Ultrafilter dichotomy and separation}
A proper filter \(U\) is an ultrafilter if and only if, for every
\(a\in B\), exactly one of \(a,\neg a\) belongs to \(U\).
Consequently,
\[
 a\vee b\in U\quad\Longleftrightarrow\quad a\in U\text{ or }b\in U.
\]
Assuming BPIT, if \(F\) is a filter and \(b\notin F\), there is an
ultrafilter \(U\supseteq F\) with \(b\notin U\).
In particular, each nonzero element belongs to an ultrafilter,
and distinct Boolean elements are separated by an ultrafilter.
The trivial Boolean algebra has no ultrafilters.
\end{hblemma}

\begin{hbdefinition}{def:p10-stone-space}{Stone space of a Boolean algebra}
A \emph{Stone space} is a compact Hausdorff topological space with
a basis of sets that are both open and closed, called \emph{clopen}
sets. The empty space is allowed.
Compact means that every open cover has a finite subcover.
For a Boolean algebra \(B\), set
\[
 X_B=\{U:U\text{ is an ultrafilter of }B\},\qquad
 \widehat a=\{U\in X_B:a\in U\}.
\]
Give \(X_B\) the topology whose opens are unions of the sets
\(\widehat a\).
The filter laws give
\(\widehat1=X_B\) and
\(\widehat a\cap\widehat b=\widehat{a\wedge b}\), so these sets
form a basis.
The theorem below proves that this is a Stone space.

For any topological space \(X\), the set
\(\operatorname{Clop}(X)\) is a Boolean algebra under intersection,
union, and complement relative to \(X\), with bounds
\(\varnothing,X\).
Thus \(\operatorname{Clop}(\varnothing)\) is the trivial Boolean
algebra. The topology on \(X_B\) is a structure on the set of
ultrafilters, distinct from the underlying set of \(B\).
\index{Stone space}\index{clopen set}\index{Stone duality}
\end{hbdefinition}

\begin{hbtheorem}{thm:p10-stone-representation}{proved=annex:p10.stone.2}{Stone representation theorem}
Assuming BPIT, \(X_B\) is a Stone space and
\[
 \beta_B:B\longrightarrow\operatorname{Clop}(X_B),
 \qquad \beta_B(a)=\widehat a,
\]
is a Boolean-algebra isomorphism.
In particular, it is an injective Boolean homomorphism onto all
clopen subsets, not merely onto a Boolean subalgebra of them.
\end{hbtheorem}

The proof establishes compactness before using it to show that a
clopen set is a finite union of basic clopens. It applies to arbitrary
open covers, without choosing a basic neighborhood independently at
every point.
The representation and the inverse construction below are
\cite[IV\S4, Theorem~4.6, pp.~136--138]{BurrisSankappanavar2012};
see also \textcite{Stone1936,Johnstone1982,GivantHalmos2009}.

\subsection{Reconstruction of spaces and action on maps}
\label{sec:p10-stone-reconstruction}

\begin{hbtheorem}{thm:p10-stone-space-reconstruction}{proved=annex:p10.stone.3}{A Stone space is recovered from its clopen algebra}
For every Stone space \(X\), the map
\[
 \varepsilon_X:X\longrightarrow X_{\operatorname{Clop}(X)},
 \qquad
 \varepsilon_X(x)=\{C\in\operatorname{Clop}(X):x\in C\},
\]
is a homeomorphism.
Every ultrafilter of \(\operatorname{Clop}(X)\) consists of exactly
the clopen neighborhoods of one point of \(X\).
For \(X=\varnothing\), both sides of the displayed map are empty.
\end{hbtheorem}

\begin{hbdefinition}{def:p10-stone-dual-maps}{The inverse-image maps of Stone duality}
Let \(\mathsf{BA}\) have Boolean algebras as objects and maps
preserving \(\wedge,\vee,\neg,0,1\) as arrows.
Let \(\mathsf{Stone}\) have Stone spaces as objects and continuous
maps as arrows. In both cases identities and composition are those
of functions.

For a Boolean homomorphism \(h:B\to C\), define
\[
 X(h):X_C\longrightarrow X_B,\qquad X(h)(U)=h^{-1}[U].
\]
For a continuous map \(g:X\to Y\), define
\[
 \operatorname{Clop}(g):
 \operatorname{Clop}(Y)\longrightarrow\operatorname{Clop}(X),
 \qquad \operatorname{Clop}(g)(D)=g^{-1}[D].
\]
Both assignments reverse the direction of arrows.
A contravariant functor preserves identities and reverses the order
of composition; the theorem verifies these laws for the displayed
assignments. Naturality of the reconstruction maps means the two
displayed commuting equations in that theorem.
These are the concrete categorical data behind the contravariant
equivalence, with the general category terminology formalized in
Definition~\ref{def:p10-category}.
\end{hbdefinition}

\begin{hbtheorem}{thm:p10-stone-duality}{proved=annex:p10.stone.4}{Functorial Stone duality}
Assuming BPIT, the inverse-image assignments are well-defined
contravariant functors
\[
 \mathsf{BA}\ \xrightarrow{\ X\ }\ \mathsf{Stone},
 \qquad
 \mathsf{Stone}\ \xrightarrow{\ \operatorname{Clop}\ }\ \mathsf{BA}.
\]
The Boolean isomorphisms \(\beta_B\) and homeomorphisms
\(\varepsilon_X\) are natural: for \(h:B\to C\) and \(g:X\to Y\),
\[
\begin{aligned}
 \operatorname{Clop}(X(h))\circ\beta_B&=\beta_C\circ h,\\
 X(\operatorname{Clop}(g))\circ\varepsilon_X&=\varepsilon_Y\circ g.
\end{aligned}
\]
They satisfy the identities
\[
\begin{aligned}
 X(\beta_B)\circ\varepsilon_{X_B}&=1_{X_B},\\
 \operatorname{Clop}(\varepsilon_X)\circ
                 \beta_{\operatorname{Clop}(X)}
      &=1_{\operatorname{Clop}(X)}.
\end{aligned}
\]
Thus the categories are contravariantly equivalent.
Every continuous map \(f:X_C\to X_B\) arises from a unique Boolean
homomorphism \(h:B\to C\) by \(f=X(h)\).
\end{hbtheorem}

The source places the action on homomorphisms and continuous maps
in exercises following the object-level representation theorem
\cite[IV\S4, Exercises~9--10, p.~140]{BurrisSankappanavar2012}.
The annex supplies the well-definedness, composition laws,
naturality calculations, and inverse map on arrows in full.

\paragraph{Finite algebras and finite spaces.}
For \(B=\mathcal P(S)\) with \(S\) finite, every ultrafilter is
\(U_s=\{A\subseteq S:s\in A\}\) for a unique \(s\in S\).
Indeed, the finite join of the singleton atoms is \(1\), so finite-join
primeness forces an ultrafilter to contain a singleton; properness
prevents it from containing two different ones.
Then upward closure and the complement criterion determine the entire
ultrafilter. Each \(\widehat{\{s\}}\) is the singleton \(\{U_s\}\),
so the dual topology is discrete.
For \(S=\varnothing\), the algebra is trivial and its dual is empty.
